% Generated draft: requires mathematical and source review.
\documentclass[11pt]{article}
\usepackage[margin=25mm]{geometry}
\usepackage{fontspec}
\setmainfont{Latin Modern Roman}
\usepackage{amsmath,amssymb,mathtools,unicode-math}
\setmathfont{Latin Modern Math}
\usepackage{enumitem,needspace,xurl,hyperref}
\setlength{\parindent}{0pt}
\setlength{\parskip}{5pt}
\setlength{\emergencystretch}{3em}
\newcommand{\N}{\mathbb{N}}
\newcommand{\Z}{\mathbb{Z}}
\newcommand{\Q}{\mathbb{Q}}
\newcommand{\R}{\mathbb{R}}
\newcommand{\C}{\mathbb{C}}
\newcommand{\F}{\mathbb{F}}
\newcommand{\A}{\mathbb{A}}
\newcommand{\PP}{\mathbb{P}}
\newcommand{\E}{\mathbb{E}}
\newcommand{\eps}{\varepsilon}
\newcommand{\dd}{\,\mathrm{d}}
\begin{document}
\section*{UnsolvedMath ID: 20002298}
\textbf{Source problem code:} catalogue entry
\subsection{Definitions and mathematical statement}
Pasten's Question \(14\) concerns the parameter-free structure
\[
 \mathcal K=(K(t);0,1,+,\cdot,T),\qquad
 T(f)\Longleftrightarrow f\notin K.
\]
Here \(T\) is a unary nonconstant predicate, not a name for the generator \(t\); the extracted ``\(f\ 6\in K\)'' means nonmembership. The source asks about existential undecidability and says the question was open for every field. Ordinary existential syntax allows negated \(T\), whereas positive existential syntax does not. The aid's reductions distinguish those readings, so the historical blanket status is not adopted as a current theorem under both conventions.

For every fixed field \(K\), its ordinary existential pure-ring theory many-one reduces to that of \(\mathcal K\): replace
\[
 \exists x_1\cdots x_m\,\theta(x_1,\ldots,x_m)
\]
by the same formula conjoined with \(\bigwedge_i\neg T(x_i)\). All witnesses are then constants, so truth is exactly truth in \(K\). In particular ordinary existential undecidability of a base field transfers. This reduction does not prove a positive-existential transfer, since the constant restriction uses a negated predicate.

If \(K=\F_q\) is finite, constants are exactly the roots of \(x^q-x\) in \(\F_q(t)\), and
\[
 \neg T(x)\Longleftrightarrow x^q=x,\qquad
 T(x)\Longleftrightarrow\exists y\ (x^q y=xy+1).
\]
Thus the positive nonconstant predicate is itself positive existentially eliminable in the pure field language. Substituting these definitions gives interreducibility of ordinary existential theories, and separately of positive existential theories, with their pure-ring counterparts. The exponent \(q\) is fixed for the field, so the substitution is effective. This clarification does not settle the reported flagship open \(\F_p(t)\) parameter-free problem.

Every fractional-linear \(K\)-automorphism
\[
 t\longmapsto\frac{at+b}{ct+d},\qquad ad-bc\ne0,
\]
preserves the predicate \(T\). The resulting \(\operatorname{PGL}_2(K)\) symmetry moves \(t\), for example by \(t\mapsto t+1\). Consequently no parameter-free formula in the expanded language defines that distinguished generator uniquely. This is an obstruction to recovering a named parameter, not to every possible undecidability interpretation. No stronger general transfer or resolution for arbitrary \(K\) is asserted.
\subsection{Short English statement}
The predicate names nonconstant functions, not the variable \(t\). The aid gives an ordinary-existential reduction from the base field, positive elimination of the predicate over finite fields, and a symmetry obstruction to defining a particular generator. These clarify two syntax readings without solving the central finite-field case.
\subsection{Sources}
\begin{itemize}
\item Local UnsolvedMath registry, permanent ID \(20002298\), source code \texttt{catalogue entry}; canonical collection \texttt{aim-logic-notes.json}, notes[73].
\item ``Problems related to `Definability and Decidability Problems in Number Theory','' AIM workshop, 9--13 September 2013, moderated by T. Scanlon and recorded by J. Demeyer, Question \(14\), Pasten: \url{https://aimath.org/pastworkshops/definabilityinntproblems.pdf}. No PDF was opened in this session.
\item Supplied research attempt \(1\), \texttt{partial\_result}, run status ``open,'' rigor label ``proved,'' novelty confidence ``low.'' Registry status \texttt{partially\_solved}, classification \texttt{PARTIAL-PROGRESS}, updated 14 August 2026. The source's ``open for every field'' is historical and syntax-sensitive.
\item No full authoritative bibliography for the reductions or current finite-field status is supplied. Original remarks and literature are empty, with no proposer-field value or year.
\end{itemize}
\end{document}
